\section{All runs}
\label{app:runs}

Table~\ref{tab:allruns} lists every run on the M5 profile together with the dense diagnostics of
Section~\ref{sec:dense}; it is generated from the committed run records by \texttt{paper/make\_tables.py}.
Table~\ref{tab:m4runs} lists every run on the earlier M4 profile, whose checkpoints are analysed in
Sections~\ref{sec:selection} and~\ref{sec:fixes}. The two profiles are never pooled. Not listed: smoke tests, and one M4 launch cancelled at about epoch 13 when its
launching machine lost its session, which the protocol classes as incomplete.

\begin{table}[h]
\centering
\scriptsize
\begin{tabular}{lrllrrrrr}
\toprule
Model & seed & numerics & data & round & held-out & best val. (epoch) & transition & epochs \\
\midrule
$S_1{=}2$ & 123 & TF32 & isolated & M5 & 0.705 & 0.704 (53) & 41 & 60 \\
dense & 123 & TF32 & isolated & M5 & 0.010 & 0.010 (55) & never & 60 \\
$S_1{=}2$ & 231 & TF32 & isolated & M5 & 0.684 & 0.685 (46) & 40 & 60 \\
$S_1{=}2$ & 312 & TF32 & isolated & M5 & 0.836 & 0.837 (46) & 35 & 60 \\
$S_1{=}12$ & 123 & TF32 & isolated & M5 & 0.850 & 0.850 (37) & 26 & 60 \\
$S_1{=}4$ & 123 & TF32 & isolated & M5 & 0.884 & 0.885 (59) & 34 & 60 \\
$S_1{=}6$ & 123 & TF32 & isolated & M5 & 0.962 & 0.962 (58) & 29 & 60 \\
dense & 810001 & TF32 & isolated & M5 & 0.012 & 0.012 (58) & never & 60 \\
$S_1{=}2$ & 810001 & TF32 & isolated & M5 & 0.578 & 0.578 (59) & 49 & 60 \\
dense & 810002 & TF32 & isolated & M5 & 1.000 & 1.000 (34) & 33 & 35 \\
$S_1{=}2$ & 810002 & TF32 & isolated & M5 & 0.709 & 0.709 (42) & 39 & 60 \\
dense & 810003 & TF32 & isolated & M5 & 0.006 & 0.007 (53) & never & 60 \\
$S_1{=}2$ & 810003 & TF32 & isolated & M5 & 0.011 & 0.011 (58) & never & 60 \\
$S_1{=}6$ & 810001 & TF32 & isolated & 2c & 0.937 & 0.937 (56) & 25 & 60 \\
$S_1{=}6$ & 810002 & TF32 & isolated & 2c & 0.945 & 0.946 (56) & 27 & 60 \\
$S_1{=}6$ & 810003 & TF32 & isolated & 2c & 0.945 & 0.944 (58) & 26 & 60 \\
dense & 123 & TF32 & legacy & check & 0.987 & 0.987 (27) & 26 & 28 \\
dense & 123 & float32 & isolated & check & 1.000 & 1.000 (42) & 41 & 43 \\
dense & 810001 & float32 & isolated & 2e & --$^\ddagger$ & 0.007 (54) & never & 55 \\
dense & 810003 & float32 & isolated & 2f & 0.007 & 0.007 (57) & never & 60 \\
\bottomrule
\end{tabular}
\caption{\textbf{Every run on the M5 profile and its diagnostics, including all non-transitions.} Epochs are zero-based; early stopping at validation recall above 0.98 ends some runs before 60 epochs. $^\ddagger$Stopped by us after 55 of 60 epochs; no held-out score.}
\label{tab:allruns}
\end{table}


\begin{table}[h]
\centering
\small
\begin{tabular}{lrrrrr}
\toprule
Model (seed) & recall & parent hit & key-block hit & recall $\mid$ hit & recall $\mid$ miss \\
\midrule
$S_1{=}2$ (123) & 0.803 & 0.888 & 0.594 & 0.999 & 0.516 \\
$S_1{=}2$ (312) & 0.715 & 0.815 & 0.489 & 0.997 & 0.446 \\
$S_1{=}12$ (123) & 0.920 & 1.000 & 0.729 & 1.000 & 0.705 \\
$S_1{=}2$ (231), never learned & 0.008 & 0.707 & 0.274 & 0.014 & 0.005 \\
\bottomrule
\end{tabular}
\caption{\textbf{Selection diagnostic on the earlier M4-profile checkpoints.} Retrieval layer, full held-out split. \emph{Parent hit}: the key's super-block was selected by at least one of the 4 heads;
\emph{key-block hit}: the key's fine block was among the $K{=}4$ blocks selected by at least one head. Recall conditioned on a key-block hit
was 0.997--1.000 for every model here that learned.}
\label{tab:selection_m4}
\end{table}

\begin{table}[h]
\centering
\scriptsize
\begin{tabular}{llrrrr}
\toprule
Experiment & Model & seed & held-out & best val.\ (epoch) & transition \\
\midrule
repair screen & $S_1{=}2$ baseline & 123 & 0.803 & 0.803 (45) & 35 \\
repair screen & $S_1{=}2$ router gate & 123 & 0.652 & 0.651 (40) & 39 \\
repair screen & $S_1{=}2$ max-pool & 123 & 0.648 & 0.648 (46) & 42 \\
repair screen & $S_1{=}2$ both & 123 & 0.751 & 0.752 (28) & 28 \\
follow-up & dense & 123 & 0.992 & 0.992 (29) & 28 \\
follow-up & $S_1{=}12$ & 123 & 0.920 & 0.920 (57) & 25 \\
follow-up & $S_1{=}2$ & 312 & 0.715 & 0.715 (58) & 40 \\
follow-up & $S_1{=}2$ & 231 & 0.008 & 0.007 (54) & never \\
\bottomrule
\end{tabular}
\caption{\textbf{Every run on the M4 profile} (float32, legacy data). All four repair-screen arms trained on
identical data from identical initial weights. Each run used its own data; seeds 231 and 312 are not paired with
seed 123.}
\label{tab:m4runs}
\end{table}

\section{Compute}
\label{app:compute}

Every run used a single NVIDIA L40S GPU, on AWS (\texttt{g6e.xlarge}) or Modal, one run per GPU at a time. M5-profile
runs took 0.9--4.6 hours each. Software versions differed slightly between machines (PyTorch 2.14.0 or 2.14.1 with
CUDA 13.0; Python 3.11 or 3.13); each run's manifest records its environment. The experiments in this paper used
roughly \$150 of cloud GPU time; the M5 selection diagnostic (inference only, twelve L40S containers on Modal) cost
under \$2.

\section{Earlier exploratory results}
\label{app:history}

Before the protocol of Section~\ref{sec:setup} was fixed, the project ran exploratory experiments that report
\emph{validation} recall from single runs with differing recipes. We list them for context only; none enters the
results above. At length 128, attention reached 0.998 against 0.097 for a gated convolution, reproducing Zoology's
recall gap. At length 2048, an unpruned selector ($S_1{=}12$) and dense attention both reached about 0.99 validation
recall, while two recurrent baselines reached 0.558 and 0.193 (the first with a different learning rate). At length
4096 every model, including dense attention, failed to learn. An early reading of these runs, that block sparsity
alone caused the recall loss at length 2048, was retracted once budget sweeps showed recall rising with the budget. These runs are not part of the public release.

\section{Reproduction}
\label{app:repro}

Code, frozen protocols, per-run manifests (seeds, SHA-256 split hashes, environment, code commit), per-epoch
metrics, training logs, the selection-diagnostic scripts (\texttt{experiments/diagnose\_routing.py},
\texttt{experiments/modal\_diagnose\_m5.py}) and their outputs, and the scripts that regenerate the figures
and the appendix tables are at
\url{https://github.com/eumnesia-labs/sparse-attention-recall} (release v1.2). Figures: \texttt{python3 paper/figures/make\_figures.py};
appendix table: \texttt{python3 paper/make\_tables.py}; each results folder's \texttt{summarize*.py} regenerates its
summary from raw records. The M5 runs used a cloud version of the training harness that also contains unrelated,
unpublished work and is not released. The release instead provides a portable runner with the same recipe
(\texttt{--task mqar --precision tf32}), and \texttt{experiments/reproduce\_m5.py} turns any stored manifest into a
local command without the cloud steps. The portable runner has not been checked for numerical equivalence with the
original runs; hardware, library versions and nondeterminism can change outcomes. Training a model requires one GPU.
Model checkpoints are not released.
